% Documentation technique du prototype de collecte d'appels à projets (PI2 n°158)
% Compilation : pdflatex documentation_code.tex (deux fois, pour la table des matières)
\documentclass[11pt,a4paper]{article}

\usepackage[utf8]{inputenc}
\usepackage[T1]{fontenc}
\usepackage[french]{babel}
\usepackage{lmodern}
\usepackage[margin=2.3cm]{geometry}
\usepackage{xcolor}
\usepackage{listings}
\usepackage{booktabs}
\usepackage{tabularx}
\usepackage{enumitem}
\usepackage{tikz}
\usetikzlibrary{positioning,arrows.meta,babel}
\usepackage[hidelinks]{hyperref}

% ---------------------------------------------------------------------------
% Couleurs et mise en forme
% ---------------------------------------------------------------------------
\definecolor{pbcolor}{RGB}{176,48,32}
\definecolor{solcolor}{RGB}{30,110,60}
\definecolor{implcolor}{RGB}{30,70,150}
\definecolor{codebg}{RGB}{246,246,246}
\definecolor{codecomment}{RGB}{110,110,110}
\definecolor{codestring}{RGB}{150,60,20}
\definecolor{codekeyword}{RGB}{20,60,150}

\lstset{
  language=Python,
  basicstyle=\ttfamily\small,
  backgroundcolor=\color{codebg},
  commentstyle=\color{codecomment}\itshape,
  stringstyle=\color{codestring},
  keywordstyle=\color{codekeyword}\bfseries,
  frame=single,
  rulecolor=\color{black!15},
  breaklines=true,
  columns=fullflexible,
  keepspaces=true,
  showstringspaces=false,
  numbers=none,
  xleftmargin=4pt,
  xrightmargin=4pt,
  aboveskip=8pt,
  belowskip=8pt,
  literate={é}{{\'e}}1 {è}{{\`e}}1 {ê}{{\^e}}1 {à}{{\`a}}1 {â}{{\^a}}1
           {ô}{{\^o}}1 {ç}{{\c c}}1 {î}{{\^i}}1 {Î}{{\^I}}1 {û}{{\^u}}1
           {ù}{{\`u}}1 {É}{{\'E}}1 {’}{{'}}1
}

% Les trois rubriques de chaque problème
\newcommand{\pb}{\par\medskip\noindent\textbf{\textcolor{pbcolor}{Problème.}}\ }
\newcommand{\sol}{\par\medskip\noindent\textbf{\textcolor{solcolor}{Solution retenue.}}\ }
\newcommand{\impl}{\par\medskip\noindent\textbf{\textcolor{implcolor}{Implémentation.}}\ }
\newcommand{\limite}{\par\medskip\noindent\textbf{Limite constatée.}\ }

\newcommand{\code}[1]{\texttt{#1}}

\setlist{itemsep=2pt,topsep=4pt}
\setlength{\parindent}{0pt}
\setlength{\parskip}{4pt}

\title{Collecte d'appels à projets -- Documentation du code\\[4pt]
\large Projet PI2 n°158 -- ESILV / DVRC}
\author{Équipe projet PI2 n°158}
\date{Novembre 2026}

\begin{document}

\maketitle
\tableofcontents
\newpage

% ===========================================================================
\section{Introduction}
% ===========================================================================

\subsection{Contexte}

Le DVRC compte environ 150 chercheurs. Chaque semaine, des financeurs publient des
appels à projets (AAP) : l'ANR, Bpifrance, la Commission européenne (Horizon Europe),
la Région Île-de-France, etc. L'objectif final du projet est d'envoyer à chaque chercheur
une newsletter ne contenant que les appels qui le concernent, après validation par
la recherche partenariale.

Le pipeline visé est le suivant :

\begin{center}
\begin{tikzpicture}[
  box/.style={draw, rounded corners, fill=black!4, minimum height=9mm,
              text width=2.1cm, align=center, font=\small},
  done/.style={box, fill=green!12, draw=green!50!black},
  arrow/.style={-{Stealth}, thick}]
  \node[done] (s) {Scraping\\hebdomadaire};
  \node[box, right=4mm of s] (db) {Base\\PostgreSQL};
  \node[box, right=4mm of db] (nlp) {NLP\\(extraction)};
  \node[box, right=4mm of nlp] (m) {Matching\\(embeddings)};
  \node[box, right=4mm of m] (e) {Éligibilité};
  \node[box, right=4mm of e] (n) {Newsletter};
  \draw[arrow] (s) -- (db);
  \draw[arrow] (db) -- (nlp);
  \draw[arrow] (nlp) -- (m);
  \draw[arrow] (m) -- (e);
  \draw[arrow] (e) -- (n);
\end{tikzpicture}
\end{center}

Le code décrit ici correspond à la première brique (en vert) : un prototype de
\emph{scraping} qui collecte les appels de quatre sources et les écrit dans un fichier CSV.

\subsection{Comment lire ce document}

Le code est organisé autour d'une série de problèmes concrets (\og comment récupérer
les dates de chaque appel ? \fg, \og comment lire un PDF scanné ? \fg\dots). Pour chacun,
le document donne :

\begin{itemize}
  \item \textbf{\textcolor{pbcolor}{le problème}} : ce qui rend la tâche difficile ;
  \item \textbf{\textcolor{solcolor}{la solution retenue}} : l'idée, indépendamment du code ;
  \item \textbf{\textcolor{implcolor}{l'implémentation}} : où et comment elle est codée,
        avec les extraits utiles.
\end{itemize}

La section~\ref{sec:transverse} traite les problèmes communs à toutes les sources, les
sections~\ref{sec:anr} à~\ref{sec:region} les problèmes propres à chaque source.
La section~\ref{sec:limites} liste les corrections déjà faites et les limites connues.

\subsection{Principes de conception}

Trois principes guident tout le code :

\begin{enumerate}
  \item \textbf{Un format commun.} Quel que soit le site, chaque appel devient une ligne
        avec exactement les mêmes colonnes. Ce qu'un site ne donne pas reste vide.
  \item \textbf{Le scraper collecte, il n'interprète pas.} Les scrapers récupèrent
        le texte brut (page web, texte des PDF, lignes de dates) et ne remplissent les champs
        structurés (date limite, statut\dots) que lorsque le site les donne clairement.
        L'interprétation fine (éligibilité, budget, date limite cachée dans un PDF) est
        laissée à l'étape NLP (section~\ref{sec:extraction}).
  \item \textbf{Code lisible par un étudiant.} Pas d'annotations de type, pas de syntaxe
        compacte, des boucles explicites et une docstring courte par fonction : le projet doit
        pouvoir être repris par une autre équipe.
\end{enumerate}

% ===========================================================================
\section{Vue d'ensemble}
% ===========================================================================

\subsection{Organisation des fichiers}

\begin{lstlisting}[language={},basicstyle=\ttfamily\scriptsize]
main.py                 lance les scrapers, écrit le CSV, affiche le taux de remplissage
config.py               paramètres (délai entre requêtes, dossiers, OCR)
scrapers/
  __init__.py           dictionnaire SCRAPERS : nom de source -> fonction scrape()
  anr.py                HTML : anr.fr/fr/appels + page de chaque appel + PDF
  bpifrance.py          HTML : bpifrance.fr/nos-appels-a-projets-concours + PDF
  horizon.py            API publique du portail Funding & Tenders (SEDIA)
  region_idf.py         open data data.iledefrance.fr (JSON)
  common.py             lecture d'une page d'appel (texte, lignes de dates, PDF)
utils/
  http_client.py        requêtes, délai, robots.txt, gestion d'erreurs
  html_tools.py         texte principal d'une page, liens PDF, lignes de dates
  pdf_tools.py          téléchargement PDF, texte (pdfplumber) puis OCR si scanné
  storage.py            format commun d'une ligne (COLUMNS) + CSV (+ PostgreSQL à venir)
  report.py             tableau de taux de remplissage par source
extraction/
  field_extractor.py    étape suivante : extraction des champs par un modèle QA (BERT)
\end{lstlisting}

Les fichiers \code{pyproject.toml}, \code{uv.lock} et \code{src/aap\_scraper/} ont été
générés par l'outil \code{uv} et ne sont pas utilisés par le prototype : les dépendances
réelles sont listées dans \code{requirements.txt}.

\subsection{Lancer le programme}

Les dépendances sont installées dans l'environnement virtuel \code{.venv}. Il faut donc
l'activer (ou appeler \code{.venv/bin/python}) avant de lancer les commandes :

Installation (une fois) :

\begin{lstlisting}[language=bash]
python -m venv .venv
source .venv/bin/activate                 # Windows : .venv\Scripts\activate
pip install -r requirements.txt
# Pour l'OCR, installer aussi Tesseract (avec les langues fra et eng) et Poppler,
# par exemple sous Arch : pacman -S tesseract tesseract-data-fra tesseract-data-eng poppler
\end{lstlisting}

Utilisation :

\begin{lstlisting}[language=bash]
source .venv/bin/activate
python main.py --max-items 3 --no-pdf     # test rapide, toutes les sources
python main.py --sources anr horizon      # seulement certaines sources
python main.py                            # tout, avec les PDF
python -m scrapers.anr                    # une seule source (idem bpifrance, horizon, region_idf)
python -m scrapers.region_idf --columns   # colonnes du jeu de données Région
\end{lstlisting}

Sorties :
\begin{itemize}
  \item \code{data/output/aap\_<date>.csv} : une ligne par appel, toutes sources confondues ;
  \item \code{data/pdfs/<source>/} : les PDF téléchargés et, à côté, leur texte en \code{.txt} ;
  \item \code{logs/run\_<date>.log} : le journal de l'exécution.
\end{itemize}

La date des noms de fichiers va jusqu'à la seconde (\code{2026-11-16\_202247}) et est
calculée une seule fois au début de \code{main()} : le CSV et le log d'une même exécution
portent le même nom, et deux exécutions rapprochées n'écrasent jamais les fichiers
l'une de l'autre.

\subsection{Déroulement d'une exécution}

\begin{enumerate}
  \item \code{main.py} configure les logs et lit les options de la ligne de commande.
  \item Pour chaque source demandée, il appelle la fonction \code{scrape()} correspondante,
        trouvée dans le dictionnaire \code{SCRAPERS} de \code{scrapers/\_\_init\_\_.py}.
  \item Chaque \code{scrape()} renvoie une liste de lignes (des dictionnaires) au format commun.
  \item \code{main.py} rassemble toutes les lignes, les écrit dans un seul CSV et affiche un
        tableau du taux de remplissage de chaque colonne par source.
\end{enumerate}

Toutes les requêtes réseau passent par \code{utils/http\_client.py}, tout le travail sur
le HTML par \code{utils/html\_tools.py} et tout le travail sur les PDF par
\code{utils/pdf\_tools.py}. Les scrapers ne contiennent donc que ce qui est propre à
leur site.

% ===========================================================================
\section{Le format commun d'un appel}
% ===========================================================================

\pb Chaque site présente ses appels différemment : une page HTML pour l'ANR, une réponse
d'API pour Horizon, un fichier JSON pour la Région. Les étapes suivantes (base de données,
NLP, matching) ne doivent pas dépendre de ces différences.

\sol Définir une fois pour toutes la liste des colonnes d'un appel. Chaque scraper part
d'une ligne vide contenant toutes ces colonnes et remplit ce qu'il peut.

\impl La liste \code{COLUMNS} et la fonction \code{new\_record()} sont dans
\code{utils/storage.py} :

\begin{lstlisting}
def new_record(source):
    """Create an empty row containing every column, so all sources share the same format."""
    record = {}
    for column in COLUMNS:
        record[column] = ""
    record["source"] = source
    record["collected_at"] = datetime.now().isoformat(timespec="seconds")
    return record
\end{lstlisting}

\begin{center}
\small
\begin{tabularx}{\linewidth}{@{}lX@{}}
\toprule
\textbf{Colonne} & \textbf{Contenu} \\
\midrule
\code{source}         & scraper qui a produit la ligne : \code{anr}, \code{bpifrance}, \code{horizon}, \code{region\_idf} \\
\code{source\_id}     & identifiant de l'appel sur son site (slug d'URL, code de topic, référence\dots) \\
\code{title}          & titre de l'appel \\
\code{funder}         & financeur (\og ANR \fg, \og Bpifrance \fg\dots) \\
\code{call\_type}     & type d'appel tel que donné par la source (\og France 2030 \fg, \og RIA \fg\dots) \\
\code{status}         & statut tel qu'écrit par la source (ouvert, \emph{forthcoming}, pré-annonce\dots) \\
\code{opening\_date}  & date d'ouverture au format \code{AAAA-MM-JJ} si elle a pu être lue \\
\code{deadline}       & date limite au format \code{AAAA-MM-JJ} si elle a pu être lue \\
\code{dates\_text}    & lignes brutes parlant de dates, gardées telles quelles pour le NLP \\
\code{summary}        & résumé court \\
\code{page\_text}     & texte principal de la page (ou description fournie par l'API) \\
\code{pdf\_urls}      & liens PDF trouvés sur la page, séparés par \og \code{ | } \fg \\
\code{pdf\_text\_files} & fichiers \code{.txt} locaux contenant le texte extrait des PDF \\
\code{pdf\_methods}   & méthode de lecture de chaque PDF : \code{text\_layer} ou \code{ocr} \\
\code{url}            & page de l'appel, pour vérifier à la main \\
\code{collected\_at}  & date et heure de la collecte \\
\bottomrule
\end{tabularx}
\end{center}

Une cellule CSV ne contenant qu'une chaîne, les listes (liens PDF, fichiers \code{.txt},
dates multiples) sont transformées en une seule chaîne par \code{join\_list()}, avec
\og \code{ | } \fg{} comme séparateur.

% ===========================================================================
\section{Problèmes communs à toutes les sources}
\label{sec:transverse}
% ===========================================================================

\subsection{Interroger les sites sans les surcharger ni enfreindre leurs règles}

\pb Un scraper envoie beaucoup de requêtes à la suite. Sans précaution, il peut surcharger
un site, ignorer les règles de son fichier \code{robots.txt}, ou s'arrêter à la première
erreur réseau.

\sol Faire passer \emph{toutes} les requêtes par un seul module qui :
\begin{itemize}
  \item attend au moins \code{REQUEST\_DELAY\_SECONDS} (1,5~s) entre deux requêtes ;
  \item lit le \code{robots.txt} de chaque site (une seule fois par site) et refuse les URL interdites ;
  \item envoie un \emph{User-Agent} explicite avec un email de contact ;
  \item transforme toute erreur (réseau, code HTTP différent de 200) en valeur \code{None}
        accompagnée d'un message dans les logs, au lieu de faire planter le programme.
\end{itemize}

\impl \code{utils/http\_client.py}. Une session \code{requests} unique est partagée par
tout le programme (connexions réutilisées, même \emph{User-Agent} partout). Le délai est
géré par une variable globale qui mémorise l'heure de la dernière requête :

\begin{lstlisting}
def wait_politely():
    """Sleep a little so that two requests are never closer than REQUEST_DELAY_SECONDS."""
    global last_request_time
    elapsed = time.time() - last_request_time
    if elapsed < config.REQUEST_DELAY_SECONDS:
        time.sleep(config.REQUEST_DELAY_SECONDS - elapsed)
    last_request_time = time.time()
\end{lstlisting}

Le \code{robots.txt} est téléchargé par \code{load\_robots()} puis analysé avec
\code{urllib.robotparser} (bibliothèque standard). Le résultat est mis en cache dans le
dictionnaire \code{robots\_cache}. Si un site n'a pas de \code{robots.txt} (code HTTP
différent de 200), tout est considéré comme autorisé.

La fonction \code{get()} enchaîne les trois étapes :

\begin{lstlisting}
def get(url, params=None, check_robots=True):
    if check_robots and not is_allowed_by_robots(url):
        log.warning("robots.txt forbids this URL, skipped: %s", url)
        return None
    wait_politely()
    try:
        response = session.get(url, params=params, timeout=config.REQUEST_TIMEOUT_SECONDS)
    except requests.RequestException as error:
        log.warning("Request failed for %s: %s", url, error)
        return None
    return check_response(response, url)
\end{lstlisting}

\code{check\_response()} renvoie la réponse seulement si le code est 200. Un code 403 donne
un message spécifique (\og access refused, maybe an anti-bot protection \fg) car c'est le
signe habituel d'une protection anti-robot.

Le paramètre \code{check\_robots=False} n'est utilisé que pour les API et fichiers open data
officiels, qui sont faits pour être téléchargés. La fonction \code{post()}, utilisée par
l'API européenne, ne vérifie pas le \code{robots.txt} pour la même raison.

\subsection{Une source en panne ne doit pas bloquer les autres}

\pb Les sites changent sans prévenir. Une erreur imprévue dans un scraper (HTML modifié,
réponse inattendue) ne doit pas faire perdre les résultats des autres sources.

\sol Lancer chaque scraper dans un bloc \code{try / except} et continuer avec une liste
vide en cas d'erreur, en gardant la trace complète de l'erreur dans les logs.

\impl Dans \code{main()} :

\begin{lstlisting}
for name in args.sources:
    log.info("===== %s =====", name)
    try:
        records = SCRAPERS[name](max_items=args.max_items, with_pdfs=not args.no_pdf)
    except Exception:
        # One broken source must not stop the others
        log.exception("The %s scraper crashed, moving on to the next source", name)
        records = []
    log.info("%s: %d calls collected", name, len(records))
    all_records.extend(records)
\end{lstlisting}

\code{log.exception()} écrit le message \emph{et} la pile d'appels, ce qui permet de
retrouver la ligne fautive dans \code{logs/}.

\subsection{Extraire le texte utile d'une page web}

\pb Une page HTML contient, en plus du texte de l'appel, des menus, un pied de page, un
bandeau de cookies, des liens de partage\dots{} Écrire un extracteur par site (en
s'appuyant sur ses classes CSS) serait long et casserait à chaque refonte du site.

\sol Utiliser \code{trafilatura}, une bibliothèque open source qui repère automatiquement
le contenu principal de n'importe quelle page. Si elle ne trouve rien, se rabattre sur tout
le texte visible de la page.

\impl \code{main\_text()} dans \code{utils/html\_tools.py} :

\begin{lstlisting}
def main_text(html):
    text = trafilatura.extract(html, include_tables=True, favor_recall=True)
    if text:
        return text
    # Fallback if trafilatura finds nothing: all the visible text of the page
    return visible_text(parse_html(html))
\end{lstlisting}

\code{include\_tables=True} garde les tableaux (souvent utilisés pour les calendriers) et
\code{favor\_recall=True} demande à trafilatura de garder plutôt trop de texte que pas
assez : il vaut mieux un peu de bruit qu'une information manquante pour l'étape NLP.

Plus généralement, \code{html\_tools.py} évite de dépendre des classes CSS des sites et
s'appuie sur des éléments stables : balise \code{<h1>}, liens finissant par \code{.pdf},
mots-clés.

\subsection{Récupérer les dates de chaque appel sur une page HTML}
\label{sec:dates-html}

\pb Les dates (ouverture, clôture, relève\dots) sont essentielles pour la newsletter, mais
chaque site les écrit à sa façon et à un endroit différent, souvent dans un encadré
latéral que l'extracteur de texte principal supprime.

\sol Deux niveaux :
\begin{enumerate}
  \item quand la source donne les dates de façon structurée (API Horizon, colonnes de la
        Région, liste des appels de l'ANR), on les lit directement et on les convertit au
        format \code{AAAA-MM-JJ} ;
  \item dans tous les cas, on garde aussi les \emph{lignes brutes} qui parlent de dates,
        repérées par mots-clés dans \emph{tout} le texte visible de la page. Elles vont
        dans \code{dates\_text} et serviront à l'étape NLP.
\end{enumerate}

\impl Les mots-clés sont dans \code{scrapers/common.py} :

\begin{lstlisting}
DATE_KEYWORDS = ["ouverture", "clôture", "cloture", "fermeture", "limite", "date", "relève", "releve",
                 "deadline"]
\end{lstlisting}

La fonction \code{short\_lines\_with\_keywords()} de \code{html\_tools.py} parcourt le texte
ligne par ligne et garde les lignes \emph{courtes} (200 caractères au plus) qui contiennent
l'un de ces mots. La limite de longueur évite de récupérer des paragraphes entiers qui
mentionnent le mot \og date \fg{} par hasard. Les doublons sont ignorés.

Sur beaucoup de pages, le libellé et la valeur sont dans deux éléments HTML différents
(\og Fermeture : \fg{} puis \og 28/01/2027 à 11:00 CET \fg), donc sur deux lignes du texte.
Une ligne qui se termine par \og : \fg{} est donc recollée à la ligne suivante, pour garder
\og Fermeture : 28/01/2027 à 11:00 CET \fg{} et pas un libellé sans valeur.

\begin{lstlisting}
def short_lines_with_keywords(text, keywords, max_length=200):
    lines = text.splitlines()
    found = []
    for index in range(len(lines)):
        line = lines[index].strip()
        if not line or len(line) > max_length:
            continue
        if line.endswith(":") and index + 1 < len(lines):
            line = line + " " + lines[index + 1].strip()
        lower_line = line.lower()
        for keyword in keywords:
            if keyword in lower_line:
                if line not in found:
                    found.append(line)
                break
    return found
\end{lstlisting}

Elle est appliquée au texte visible complet (\code{visible\_text()}), et non au texte
principal de trafilatura, justement parce que les dates sont souvent dans un bloc latéral.

\subsection{Mettre toutes les dates au même format}

\pb Les dates arrivent sous plusieurs formes : \code{16/11/2026} (sites français),
\code{2026-11-16T17:00:00.000+0000} (API européenne), \code{2026-08-31T22:00:00+00:00}
(open data Région). Pour trier et filtrer les appels, il faut un format unique.

Piège : une date avec fuseau horaire ne peut pas être coupée telle quelle.
\code{2026-08-31T22:00:00+00:00} est en temps universel (UTC) ; à Paris, c'est le
1\textsuperscript{er} septembre à minuit. Garder les dix premiers caractères donnerait
le 31 août, soit un jour d'écart.

\sol Convertir toutes les dates au format ISO \code{AAAA-MM-JJ}, qui se trie correctement
comme du texte et que PostgreSQL comprend directement. Une date qui porte un fuseau
horaire est d'abord convertie à l'heure de Paris (tous les appels collectés sont français
ou européens).

\impl \code{to\_iso\_date()} dans \code{utils/storage.py} essaie d'abord de lire une
date-heure complète avec fuseau (\code{to\_paris\_date()}), puis cherche une date
française avec une expression régulière, puis une date ISO, et renvoie une chaîne vide
si elle ne trouve rien :

\begin{lstlisting}
PARIS = ZoneInfo("Europe/Paris")

def to_paris_date(text):
    """Read a full ISO date-time with a time zone and return its day in Paris ('' if not possible)."""
    try:
        moment = datetime.fromisoformat(text.strip())
    except ValueError:
        return ""
    if moment.tzinfo is None:
        return ""  # no time zone given: to_iso_date() keeps the date as written
    return moment.astimezone(PARIS).date().isoformat()

DATE_FR = re.compile(r"(\d{2})/(\d{2})/(\d{4})")
DATE_ISO = re.compile(r"\d{4}-\d{2}-\d{2}")

def to_iso_date(text):
    if not text:
        return ""
    paris_date = to_paris_date(text)
    if paris_date:
        return paris_date
    match = DATE_FR.search(text)
    if match:
        day, month, year = match.groups()
        return year + "-" + month + "-" + day
    match = DATE_ISO.search(text)
    if match:
        return match.group(0)
    return ""
\end{lstlisting}

\code{ZoneInfo} vient de la bibliothèque standard (module \code{zoneinfo}) ; sous Windows,
qui n'a pas de base de fuseaux horaires, il faut le petit paquet \code{tzdata}
(présent dans \code{requirements.txt}). La fonction \code{today\_iso()} donne la date du
jour à Paris au même format, pour comparer les dates des appels à aujourd'hui.

\subsection{Trouver les documents PDF d'un appel}

\pb Le texte complet d'un appel (critères d'éligibilité, budget, calendrier détaillé) est
souvent uniquement dans un PDF : texte de l'appel, cahier des charges, annexes.

\sol Garder tous les liens de la page dont l'adresse contient \code{.pdf}, les rendre
absolus et supprimer les doublons.

\impl \code{find\_pdf\_links()} dans \code{html\_tools.py}. \code{urljoin()} transforme
un lien relatif (\code{/medias/aap.pdf}) en adresse complète à partir de l'URL de la page.

Point d'attention dans \code{scrapers/common.py} : les liens PDF sont cherchés
\emph{avant} l'appel à \code{visible\_text()}, car cette dernière supprime des balises de
l'objet \code{soup} (\code{script}, \code{style}, \code{noscript}).

\subsection{Lire le texte d'un PDF, y compris s'il est scanné}

\pb Il existe deux sortes de PDF : les PDF \og numériques \fg, qui contiennent le texte, et
les PDF \og image \fg{} (scans, ou documents dont le texte a été transformé en images), qui ne
contiennent que des images. Pour les seconds, il faut de la reconnaissance de caractères
(OCR), qui est lente et demande des logiciels installés sur l'ordinateur.

\sol
\begin{enumerate}
  \item Lire la couche texte avec \code{pdfplumber} (rapide et exact).
  \item Calculer le nombre moyen de caractères par page. S'il est inférieur à
        \code{OCR\_MIN\_CHARS\_PER\_PAGE} (100), le PDF est considéré comme scanné.
  \item Seulement dans ce cas, essayer l'OCR avec Tesseract (open source, tourne en local),
        en français et en anglais.
  \item Garder le texte le plus long des deux.
\end{enumerate}

\impl \code{utils/pdf\_tools.py}.

\code{download\_pdf()} enregistre le fichier dans \code{data/pdfs/<source>/}. Si le fichier
existe déjà (exécution précédente), il n'est pas retéléchargé. Comme certains liens
\og .pdf \fg{} mènent en réalité à une page HTML, le contenu est vérifié : un vrai PDF
commence toujours par les octets \code{\%PDF}.

\code{pdf\_to\_text()} applique la stratégie ci-dessus :

\begin{lstlisting}
chars_per_page = len(text.strip()) / max(nb_pages, 1)
if chars_per_page >= config.OCR_MIN_CHARS_PER_PAGE:
    return text, "text_layer"

log.info("Little text found in %s, trying OCR", path)
ocr_text = extract_text_with_ocr(path)
if len(ocr_text.strip()) > len(text.strip()):
    return ocr_text, "ocr"
return text, "text_layer"
\end{lstlisting}

L'OCR est indispensable en pratique : les textes des appels France 2030 publiés par l'ANR
sont des PDF dont le texte a été transformé en images (environ 0 caractère par page pour
pdfplumber). Il reste toutefois \emph{optionnel} pour le programme :
\code{extract\_text\_with\_ocr()} importe \code{pytesseract} et \code{pdf2image} à
l'intérieur de la fonction, dans un \code{try / except ImportError}. Si ces bibliothèques
(ou les logiciels Tesseract et Poppler) ne sont pas installées, le programme continue avec
un simple message dans les logs, et le texte de ces PDF reste vide.

\code{process\_pdfs()} traite au plus \code{MAX\_PDFS\_PER\_CALL} (3) PDF par appel : les
premiers liens d'une page sont en général le texte de l'appel. Le texte de chaque PDF est
écrit dans un fichier \code{.txt} à côté du PDF, et c'est le chemin de ce fichier qui est
mis dans le CSV (le texte complet d'un PDF serait trop long pour une cellule).

\subsection{Donner un nom unique à chaque PDF}

\pb Deux appels différents peuvent avoir un PDF portant le même nom
(\code{annexe.pdf}), et certains noms sont très longs.

\sol Préfixer le nom par une empreinte de l'URL complète et ne garder que la fin du nom
d'origine, pour rester lisible.

\impl \code{pdf\_file\_name()} : les 10 premiers caractères du hachage MD5 de l'URL, un
\og \_ \fg, puis les 60 derniers caractères du nom d'origine (avec \code{.pdf} ajouté s'il
manque). Exemple : \code{2371e5e895\_france2030-aap-labex-2030-2026.pdf}.

\subsection{Enregistrer les résultats dans un format exploitable}

\pb Les résultats doivent pouvoir être ouverts dans Excel par l'encadrant, et plus tard
chargés dans PostgreSQL.

\sol Un fichier CSV par exécution, encodé en \code{utf-8-sig} : la signature au début du
fichier indique à Excel qu'il s'agit d'UTF-8, sans quoi les accents s'affichent mal.

\impl \code{save\_to\_csv()} dans \code{utils/storage.py} utilise \code{csv.DictWriter}
avec la liste \code{COLUMNS} comme en-tête. Excel n'affiche pas plus de 32\,767 caractères
par cellule : les textes plus longs que \code{MAX\_CELL\_CHARS} (30\,000) sont coupés et
terminés par \og \code{[...]} \fg. Ce n'est pas une perte, car le texte complet des PDF est
dans les fichiers \code{.txt}.

\code{save\_to\_postgres()} est un emplacement réservé, qui n'écrit rien pour l'instant. Sa
docstring propose d'utiliser le couple (\code{source}, \code{source\_id}) comme clé, pour ne
pas créer de doublons d'une semaine à l'autre.

\subsection{Savoir ce que chaque source fournit vraiment}

\pb Toutes les sources ne donnent pas les mêmes informations. Il faut voir rapidement,
après chaque exécution, quels champs sont remplis pour quelle source, afin de savoir ce
que l'étape NLP devra chercher dans les textes.

\sol Afficher à la fin de chaque exécution un tableau : une ligne par source, une colonne
par champ, avec le pourcentage de lignes où ce champ est rempli.

\impl \code{utils/report.py}. \code{list\_sources()} liste les sources dans l'ordre
d'apparition, \code{fill\_rate()} calcule le pourcentage de lignes non vides pour une
colonne, et \code{print\_fill\_report()} affiche le tableau pour les colonnes de
\code{REPORT\_COLUMNS}. Exemple de sortie :

\begin{lstlisting}[language={}]
source        rows   title  opening_date  deadline  dates_text  summary ...
horizon          3    100%          100%      100%        100%       0% ...
\end{lstlisting}

\subsection{Garder une trace de chaque exécution}

\pb Le scraping sera lancé chaque semaine, souvent sans personne devant l'écran. Il faut
pouvoir comprendre après coup ce qui s'est passé.

\sol Écrire les messages à la fois dans le terminal et dans un fichier de log daté.

\impl \code{setup\_logging()} dans \code{main.py} crée \code{logs/run\_<date>.log} et
configure deux sorties (terminal et fichier). Les bibliothèques \code{trafilatura} et
\code{pdfminer} (utilisée par pdfplumber) étant très bavardes, seuls leurs avertissements
sont gardés.

\subsection{Modifier les réglages sans toucher au code}

\pb Délai entre requêtes, dossiers de sortie, seuil d'OCR\dots{} sont des valeurs qu'on
veut pouvoir ajuster sans chercher dans tout le code.

\sol Les regrouper dans un seul fichier.

\impl \code{config.py} : \code{USER\_AGENT}, \code{REQUEST\_DELAY\_SECONDS},
\code{REQUEST\_TIMEOUT\_SECONDS}, \code{OUTPUT\_DIR}, \code{PDF\_DIR}, \code{LOG\_DIR},
\code{MAX\_PDFS\_PER\_CALL}, \code{OCR\_MIN\_CHARS\_PER\_PAGE}, \code{OCR\_LANGUAGES}. Avant
la mise en production, il faut remplacer l'email générique du \code{USER\_AGENT} par
une vraie adresse de contact.

\subsection{Lire la page de détail d'un appel (sources HTML)}

\pb L'ANR et Bpifrance fonctionnent de la même manière : une page liste les appels, puis
chaque appel a sa propre page avec le texte complet, les dates et les PDF. Le travail sur
la page de détail est le même pour les deux sites.

\sol Une fonction commune qui, à partir de l'URL d'un appel, remplit le texte, les lignes
de dates, les liens PDF et, si demandé, le texte des PDF. Elle renvoie la page analysée,
pour qu'un scraper puisse y lire des informations propres à son site sans la télécharger
une seconde fois (c'est ce que fait l'ANR pour son calendrier).

\impl \code{add\_page\_details()} dans \code{scrapers/common.py} :

\begin{lstlisting}
def add_page_details(record, source, with_pdfs):
    response = http_client.get(record["url"])
    if response is None:
        return None

    soup = html_tools.parse_html(response.text)
    if not record["title"]:
        record["title"] = html_tools.page_title(soup)

    # PDF links first (visible_text() below removes some tags from the soup)
    pdf_urls = html_tools.find_pdf_links(soup, record["url"])
    record["pdf_urls"] = join_list(pdf_urls)

    record["page_text"] = html_tools.main_text(response.text)

    full_text = html_tools.visible_text(soup)
    date_lines = html_tools.short_lines_with_keywords(full_text, DATE_KEYWORDS)
    record["dates_text"] = join_list(date_lines)

    if with_pdfs and pdf_urls:
        txt_files, methods = pdf_tools.process_pdfs(pdf_urls, source)
        record["pdf_text_files"] = join_list(txt_files)
        record["pdf_methods"] = join_list(methods)
    return soup
\end{lstlisting}

Si le titre n'a pas été trouvé sur la page de liste, il est pris dans la page de détail
(\code{page\_title()} : la balise \code{<h1>}, sinon la balise \code{<title>}).

% ===========================================================================
\section{ANR (Agence nationale de la recherche)}
\label{sec:anr}
% ===========================================================================

Fichier : \code{scrapers/anr.py}. Méthode : lecture du HTML. Page de liste :
\url{https://anr.fr/fr/appels/}.

\subsection{Lister tous les appels en cours et à venir}

\pb Il faut trouver, sur la page de liste, chaque appel et le lien vers sa page.

\sol Sur cette page, tous les appels sont affichés d'un coup (pas de pagination) et chaque
appel a un titre \code{<h2>} contenant un lien vers une adresse qui contient
\code{/call/}. Ce motif est stable et ne dépend pas des classes CSS.

\impl \code{get\_call\_cards()} parcourt tous les \code{<h2>} et garde ceux dont le lien
contient \code{/call/}. Pour chacun, il garde le titre, l'URL absolue et le texte du bloc
de l'appel (voir ci-dessous).

\subsection{Isoler le bloc HTML d'un seul appel}

\pb Les informations d'un appel (type, période, résumé, mention \og Pré-annonce \fg) sont
autour de son titre, dans un bloc (\og carte \fg) dont on ne veut pas dépendre de la classe CSS.

\sol Partir du \code{<h2>} et remonter dans l'arbre HTML tant que le bloc parent ne contient
qu'un seul \code{<h2>}. Dès qu'un parent en contient plusieurs, c'est qu'il regroupe
plusieurs appels : on s'arrête juste avant.

\impl \code{find\_card\_block()} (4 niveaux au plus, par sécurité) :

\begin{lstlisting}
def find_card_block(h2):
    block = h2
    for _ in range(4):
        parent = block.parent
        if parent is None:
            break
        if len(parent.find_all("h2")) > 1:
            break
        block = parent
    return block
\end{lstlisting}

\subsection{Récupérer le type, le statut et le résumé}

\pb Ces informations sont dans le texte de la carte, sans balise dédiée.

\sol Exploiter la mise en page constante des cartes :
\begin{itemize}
  \item la première ligne de la carte est le type (\og Appel à projets spécifique \fg, \og France 2030 \fg) ;
  \item si le texte contient \og pré-annonce \fg, le statut est \og pré-annonce \fg, sinon \og ouvert \fg ;
  \item le résumé est la ligne la plus longue de la carte, en excluant le titre et le lien
        \og Plus de détails \fg{} (les autres lignes sont des étiquettes courtes).
\end{itemize}

\impl \code{card\_to\_record()} et \code{read\_summary()}.

\subsection{Récupérer les dates de chaque appel}

\pb Les dates apparaissent à deux endroits :
\begin{itemize}
  \item sur la page de liste, chaque carte affiche une période, soit avec des dates exactes
        (\og 10/11/2026 - 28/01/2027 \fg), soit avec des mois seulement pour les pré-annonces
        (\og Décembre 2026 - Février 2027 \fg) ;
  \item sur la page de l'appel, un calendrier donne chaque étape (\og Ouverture : \fg,
        \og Fermeture : \fg\dots) avec l'heure.
\end{itemize}
Deux pièges : dans le HTML de la carte, les deux dates de la période sont séparées par un
retour à la ligne (une recherche ligne par ligne ne trouve donc jamais la plage complète),
et pour une pré-annonce, le calendrier contient une date \emph{estimée} (\og Février 2027 \fg{}
affiché, mais \code{2027-02-15} dans le code HTML).

\sol
\begin{enumerate}
  \item Lire la période dans le bloc de dates de la carte, en remplaçant tous les retours à
        la ligne par des espaces. Si elle contient deux dates exactes, la première est
        l'ouverture et la seconde la date limite.
  \item Sur la page de l'appel, lire le calendrier : chaque étape est un bloc
        \code{<div class="task">} dont l'attribut \code{data-date} donne la date au format
        ISO. Ces dates, plus précises, remplacent celles de la liste. On ne garde que les
        étapes dont le texte affiché est une vraie date \code{jj/mm/aaaa}, pour ne jamais
        présenter une estimation comme une date exacte.
  \item Garder dans \code{dates\_text} la période de la liste suivie des lignes de dates de
        la page, pour l'étape NLP.
\end{enumerate}

\impl \code{read\_period()} lit le bloc \code{class="date"} de la carte :

\begin{lstlisting}
def read_period(block):
    date_block = block.find(class_="date")
    if date_block is None:
        return ""
    words = date_block.get_text(" ").split()
    return " ".join(words)
\end{lstlisting}

\code{card\_to\_record()} y cherche ensuite la plage avec l'expression régulière
\code{DATE\_RANGE}. Le calendrier de la page est lu par \code{read\_calendar()}, qui
renvoie une liste de couples (libellé, date ISO) :

\begin{lstlisting}
EXACT_DATE = re.compile(r"\d{2}/\d{2}/\d{4}")

def read_calendar(soup):
    steps = []
    for task in soup.find_all("div", class_="task"):
        if not EXACT_DATE.search(task.get_text(" ")):
            continue
        title = task.find(class_="task-title")
        if title is None:
            label = ""
        else:
            label = title.get_text(" ", strip=True)
        steps.append((label, task.get("data-date", "")))
    return steps
\end{lstlisting}

\code{add\_calendar\_dates()} parcourt ces étapes : la première dont le libellé contient
un mot de \code{OPENING\_LABELS} (\og ouverture \fg) donne \code{opening\_date}, la
première dont le libellé contient un mot de \code{DEADLINE\_LABELS} (\og fermeture \fg,
\og clôture \fg, \og limite \fg) donne \code{deadline}. Pour un appel en deux étapes
(pré-propositions puis propositions complètes), c'est donc la première date limite qui
est gardée : c'est celle à ne pas manquer.

Dans \code{scrape()}, la page renvoyée par \code{add\_page\_details()} est réutilisée :

\begin{lstlisting}
record = card_to_record(card)
listing_dates = record["dates_text"]
soup = add_page_details(record, SOURCE, with_pdfs)
if soup is not None:
    add_calendar_dates(record, soup)
# Keep the period of the listing in front of the detailed date lines
if listing_dates:
    record["dates_text"] = listing_dates + " | " + record["dates_text"]
\end{lstlisting}

Résultat au 16/11/2026 : 18 appels sur 24 ont des dates exactes ; les 6 autres sont des
pré-annonces, qui n'affichent que des mois (gardés dans \code{dates\_text}).

\subsection{Identifier un appel de façon stable}

\pb Il faut un identifiant pour reconnaître le même appel d'une semaine à l'autre.

\sol Le dernier segment de l'URL (le \emph{slug}) est unique et stable.

\impl \code{record["source\_id"] = card["url"].rstrip("/").split("/")[-1]}.

\subsection{Texte complet et PDF}

Gérés par \code{add\_page\_details()} (section~\ref{sec:transverse}). Les PDF de l'ANR
sont le texte de l'appel, les annexes, les FAQ et les supports de webinaire.

% ===========================================================================
\section{Bpifrance}
\label{sec:bpifrance}
% ===========================================================================

Fichier : \code{scrapers/bpifrance.py}. Méthode : lecture du HTML (Bpifrance n'a pas
d'API). Page de liste : \url{https://www.bpifrance.fr/nos-appels-a-projets-concours}.

\subsection{Lister les appels répartis sur plusieurs pages}

\pb La liste est découpée en pages (\code{?page=0}, \code{?page=1}\dots) et on ne connaît
pas le nombre de pages à l'avance.

\sol Lire les pages une par une et s'arrêter dès qu'une page n'apporte aucun nouveau lien
d'appel. Une limite de 15 pages (\code{MAX\_LIST\_PAGES}) évite une boucle infinie si la
pagination ne s'arrête jamais.

\impl \code{collect\_call\_links()}. Un lien est retenu s'il contient
\code{/nos-appels-a-projets-concours/} et n'est pas la page de liste elle-même.
\code{clean\_url()} rend le lien absolu et supprime ce qui suit \code{?} ou \code{\#}, pour
que le même appel ne soit pas compté deux fois. La boucle s'arrête aussi dès que
\code{max\_items} liens ont été trouvés.

\subsection{Remplir les champs de chaque appel}

\pb Bpifrance opère surtout des appels pour le compte de l'État (France 2030\dots) et
n'affiche pas le type ni le statut de façon uniforme.

\sol Ne remplir que ce qui est fiable (financeur, URL, identifiant, titre), et laisser
type et statut vides : ils seront trouvés par l'étape d'extraction.

\impl Dans \code{scrape()}, \code{source\_id} est la fin de l'URL. Tout le reste (titre
depuis le \code{<h1>}, texte, lignes de dates, PDF des cahiers des charges) vient de
\code{add\_page\_details()}.

\subsection{Détecter un refus d'accès}

\pb Le site peut refuser les robots (protection anti-bot).

\sol Ne pas tenter de contourner la protection, mais le signaler clairement.

\impl Le code HTTP 403 est détecté par \code{check\_response()} dans \code{http\_client.py},
qui écrit un message explicite. \code{collect\_call\_links()} s'arrête alors et le scraper
renvoie une liste vide. Les conditions générales d'utilisation du site doivent aussi être
vérifiées à la main avant une utilisation hebdomadaire.

\limite C'est ce qui se produit actuellement : tout le site, \code{robots.txt} compris,
renvoie une erreur 403 générée par le CDN CloudFront, avant même d'atteindre le site.
Contourner cette protection (par exemple en se faisant passer pour un navigateur) serait
contraire aux règles du projet. Les pistes sont de demander une autorisation à Bpifrance
ou de trouver une source qui reprend ses appels. Voir section~\ref{sec:limites}.

% ===========================================================================
\section{Horizon Europe}
\label{sec:horizon}
% ===========================================================================

Fichier : \code{scrapers/horizon.py}. Méthode : API publique du portail
\emph{EU Funding \& Tenders}.

\subsection{Récupérer les données d'un site entièrement en JavaScript}

\pb Le portail européen est une application JavaScript : le HTML téléchargé ne contient
pas les appels, ils sont chargés ensuite par le navigateur.

\sol Interroger directement l'API de recherche publique que le portail utilise lui-même.
Elle ne demande pas de compte : la clé \code{apiKey=SEDIA} est la clé publique du portail.

\impl \code{fetch\_page()}. L'API attend une requête POST où les filtres sont envoyés comme
des \og fichiers \fg{} JSON (format \emph{multipart}), d'où l'usage du paramètre
\code{files} de \code{requests} :

\begin{lstlisting}
params = {"apiKey": "SEDIA", "text": "***", "pageSize": PAGE_SIZE, "pageNumber": page_number}
files = {
    "query": ("blob", json.dumps(build_query()), "application/json"),
    "languages": ("blob", json.dumps(["en"]), "application/json"),
    "sort": ("blob", json.dumps({"order": "ASC", "field": "identifier"}), "application/json"),
}
response = http_client.post(API_URL, params=params, files=files)
\end{lstlisting}

\subsection{Ne garder que les appels Horizon Europe ouverts ou à venir}

\pb Le portail contient tous les programmes européens et tous les appels, y compris clos.

\sol Filtrer dans la requête avec les codes numériques du portail (peu documentés,
retrouvés en observant les requêtes du site) :
\begin{itemize}
  \item \code{frameworkProgramme} = \code{43108390} : Horizon Europe ;
  \item \code{status} = \code{31094501} (\emph{forthcoming}) ou \code{31094502} (\emph{open}) ;
  \item \code{type} = \code{1} ou \code{2} : topics et appels de subventions.
\end{itemize}

\impl \code{build\_query()} renvoie un filtre de type Elasticsearch (\code{bool} /
\code{must} / \code{terms}) contenant ces trois conditions.

\subsection{Écarter les entrées périmées de l'index}

\pb L'index de l'API contient des entrées anciennes (indexées en 2022--2023 et jamais mises
à jour) qui ont toujours le statut \emph{open} ou \emph{forthcoming} alors que toutes leurs
dates limites sont passées. Au 16/11/2026, c'était le cas de 458 topics sur 787. Ce sont
aussi celles qui n'ont pas de description. Un filtre sur la date ajouté à la requête est
ignoré par l'API.

\sol Faire le tri nous-mêmes : un topic dont la dernière date limite est passée est ignoré.
Un topic sans date limite est gardé (l'étape NLP lira son texte).

\impl \code{is\_outdated()}, appelée dans \code{scrape()} pour chaque résultat. Le nombre de
topics écartés est écrit dans les logs. Les dates étant des chaînes \code{AAAA-MM-JJ},
elles se comparent directement avec \code{today\_iso()} :

\begin{lstlisting}
def is_outdated(record):
    """Tell if a topic is an old entry of the index: its (last) deadline is already past."""
    if not record["deadline"]:
        return False  # no deadline given: we keep it, the NLP step will look at the text
    return record["deadline"] < today_iso()
\end{lstlisting}

\subsection{Parcourir toutes les pages de résultats sans en perdre}

\pb L'API renvoie les résultats par pages de 50. Si l'ordre des résultats n'est pas
strictement défini, il peut changer d'une page à l'autre : c'était le cas avec un tri par
date limite, car beaucoup de topics ont la même date. Résultat observé : 787 topics
annoncés, mais seulement 771 différents reçus, certains en double et d'autres jamais.

\sol Trier sur un champ unique, l'identifiant du topic, puis demander les pages une par une
jusqu'à recevoir une page incomplète (moins de 50 résultats), avec une limite de sécurité
de 40 pages (2\,000 topics).

\impl Champ \code{sort} de \code{fetch\_page()} et boucle \code{while} dans
\code{scrape()}. Le nombre total annoncé par l'API est affiché dans les logs à la première
page. Si \code{max\_items} est atteint, la fonction s'arrête immédiatement.

\subsection{Éviter les doublons}

\pb Par sécurité, on veut être sûr de ne jamais avoir deux fois le même topic.

\sol Mémoriser les identifiants déjà vus et ignorer les répétitions.

\impl Liste \code{seen\_ids} dans \code{scrape()} ; les résultats sans identifiant sont
aussi ignorés.

\subsection{Transformer un résultat de l'API en ligne du format commun}

\pb Dans la réponse de l'API, chaque champ est une liste (\code{"title": ["..."]}), même
quand il n'a qu'une valeur, et les statuts sont des codes numériques.

\sol Une petite fonction qui renvoie le premier élément de la liste (ou une chaîne vide),
et un dictionnaire de correspondance pour les statuts.

\impl \code{first\_value()} et \code{result\_to\_record()} :
\begin{itemize}
  \item \code{status} : traduit par \code{STATUS\_LABELS} (\emph{forthcoming}, \emph{open},
        \emph{closed}), ou code brut si inconnu ;
  \item \code{call\_type} : champ \code{typesOfAction} (ex. \og HORIZON Research and Innovation
        Actions \fg), avec les espaces doublés supprimés ;
  \item \code{opening\_date} : champ \code{startDate} ;
  \item \code{deadline} : certains topics ont plusieurs dates limites (topics en deux étapes,
        dates de relève de l'EIC). Elles sont triées, \code{next\_deadline()} garde la
        prochaine date qui n'est pas encore passée (ou la dernière si toutes le sont), et
        toutes sont listées dans \code{dates\_text} ;
  \item \code{summary} : titre et identifiant de l'appel parent (\code{callTitle},
        \code{callIdentifier}) ;
  \item \code{page\_text} : la description du topic (\code{descriptionByte}, en HTML),
        convertie en texte par \code{html\_to\_text()} ;
  \item \code{url} : construite à partir de l'identifiant du topic, vers la page publique du portail.
\end{itemize}

\subsection{Pourquoi pas de PDF pour cette source}

La description du topic est déjà dans la réponse de l'API. Les PDF associés sont les
\emph{work programmes}, des documents de plusieurs centaines de pages communs à des
centaines de topics : les télécharger n'apporterait presque rien. Le paramètre
\code{with\_pdfs} est accepté pour garder la même signature que les autres scrapers,
mais il n'est pas utilisé.

% ===========================================================================
\section{Région Île-de-France}
\label{sec:region}
% ===========================================================================

Fichier : \code{scrapers/region\_idf.py}. Méthode : jeu de données open data
\og Aides et appels à projets de la Région \fg{} (celui qui alimente
\url{https://www.iledefrance.fr/aides-et-appels-a-projets}).

\subsection{Récupérer toutes les aides sans scraper le site}

\pb Le site de la Région est construit à partir d'une base de données. Scraper ses pages
serait plus lent et plus fragile que lire la base directement.

\sol Télécharger en une seule requête le jeu de données complet au format JSON depuis le
portail open data (Opendatasoft). En cas d'échec, utiliser la copie publiée sur
data.gouv.fr.

\impl \code{download\_rows()} essaie les deux URL (\code{DATA\_URL} puis
\code{BACKUP\_DATA\_URL}) et renvoie la première liste valide. Le \code{robots.txt} n'est
pas vérifié (\code{check\_robots=False}) car il s'agit d'un fichier open data fait pour être
téléchargé.

\subsection{Trouver les bonnes colonnes}

\pb Les noms des colonnes du jeu de données sont peu documentés, et deux colonnes peuvent
porter la même information (\code{titre\_alternatif\_de\_l\_aide}, rempli pour 144 aides
seulement, et \code{nom\_de\_l\_aide\_de\_la\_demarche}, toujours rempli mais souvent en
majuscules).

\sol Relever une fois les vrais noms de colonnes et les écrire en constantes en haut du
fichier. Quand plusieurs colonnes peuvent donner un champ, les essayer dans l'ordre de
préférence et garder la première non vide. Une commande affiche les colonnes du jeu de
données, pour mettre à jour les constantes si la Région les renomme.

\impl Constantes \code{ID\_COLUMN}, \code{TITLE\_COLUMNS}, \code{SUMMARY\_COLUMNS},
\code{OPENING\_COLUMN}, \code{DEADLINE\_COLUMN}, \code{NEXT\_CAMPAIGN\_COLUMN},
\code{URL\_COLUMN} (vérifiées le 16/11/2026), et la fonction :

\begin{lstlisting}
def first_filled(row, columns):
    """Return the value of the first of these columns that is not empty ('' if none)."""
    for column_name in columns:
        value = to_text(row.get(column_name))
        if value:
            return value
    return ""
\end{lstlisting}

\code{python -m scrapers.region\_idf -{}-columns} appelle \code{print\_columns()}, qui affiche
le nom de chaque colonne avec le début de sa valeur pour la première aide.

\subsection{Gérer des valeurs de types variés}

\pb Selon la colonne, une valeur peut être du texte, un nombre, une liste ou \code{None}.

\sol Tout convertir en texte simple avant usage.

\impl \code{to\_text()} : \code{None} devient \code{""}, une liste devient ses éléments
séparés par des virgules, le reste passe par \code{str()}. Les colonnes contenant du HTML
(titre, résumé) sont nettoyées par \code{html\_to\_text()}.

\subsection{Récupérer les dates}

\pb Les dates (\code{date\_ouverture}, \code{date\_cloture}) sont en temps universel :
\code{2026-08-31T22:00:00+00:00} correspond au 1\textsuperscript{er} septembre à minuit à
Paris. Une troisième colonne, \code{datedebutfuturcampagne}, annonce parfois la prochaine
campagne d'une aide.

\sol Convertir les dates à l'heure de Paris (c'est ce que fait \code{to\_iso\_date()},
voir section~\ref{sec:transverse}). La date de prochaine campagne n'est ni une ouverture
ni une clôture : elle est seulement ajoutée à \code{dates\_text}.

\impl Dans \code{row\_to\_record()} : \code{opening\_date} et \code{deadline} sont remplies
avec \code{to\_iso\_date()}, et \code{dates\_text} contient par exemple
\og \code{ouverture: 2026-09-01 | clôture: 2027-06-30 | prochaine campagne: 2027-03-03} \fg.

\subsection{Donner un statut à chaque aide}

\pb Le jeu de données n'a aucune colonne de statut, et il contient aussi les aides closes
(146 sur 357 au 16/11/2026).

\sol Calculer le statut à partir des dates et de la date du jour. Les aides closes sont
gardées, avec le statut \og clos \fg, pour que le choix de les ignorer reste visible et se
fasse à l'étape suivante.

\impl \code{compute\_status()} :

\begin{lstlisting}
def compute_status(opening_date, deadline):
    today = today_iso()
    if deadline and deadline < today:
        return "clos"
    if opening_date and opening_date > today:
        return "à venir"
    if opening_date or deadline:
        return "ouvert"
    return "sans date"
\end{lstlisting}

Il n'y a pas non plus de colonne de type d'appel, et les titres ne permettent pas de le
deviner de façon fiable (près de 90~\% ne contiennent ni \og appel à projets \fg, ni
\og AMI \fg, ni mot équivalent) : \code{call\_type} reste vide pour cette source.

\subsection{Conserver toute l'information utile pour l'étape NLP}

\pb Les informations utiles (publics éligibles, objectifs, modalités) sont réparties sur
une trentaine de colonnes. Beaucoup existent en deux versions (\code{objectif} en HTML et
\code{objectif\_txt} en texte), et certaines sont purement techniques (identifiants,
codes de thèmes, SIREN\dots).

\sol Rassembler dans \code{page\_text} les colonnes utiles non vides, sous la forme
\og nom\_de\_colonne : valeur \fg, une par ligne, sans les doublons ni les colonnes
techniques.

\impl \code{row\_to\_full\_text()} saute les colonnes de \code{TECHNICAL\_COLUMNS} et toute
colonne \code{X} pour laquelle une colonne \code{X\_txt} existe.

\subsection{Récupérer les documents PDF de chaque aide}

\pb Les documents d'une aide ne sont pas dans le même jeu de données, et ce ne sont pas
tous des PDF (Word, Excel, archives ZIP\dots).

\sol Télécharger le second jeu de données, \code{documents-publics-des-aides-regionales},
et le relier aux aides par la référence administrative. Ne garder que les PDF, en mettant
les \og documents d'information \fg{} avant les \og documents nécessaires \fg{} (surtout des
formulaires à remplir), car seuls les premiers PDF d'un appel sont lus
(\code{MAX\_PDFS\_PER\_CALL}).

\impl \code{download\_pdf\_links()} renvoie un dictionnaire
\{référence administrative : liste d'URL de PDF\}. Dans \code{scrape()}, ces URL vont dans
\code{pdf\_urls} et, si \code{with\_pdfs} est vrai, sont traitées par
\code{pdf\_tools.process\_pdfs()} comme pour les autres sources. Au 16/11/2026, 226 aides
sur 357 ont au moins un PDF.

\subsection{Pas de filtrage par thème à la collecte}

Le jeu de données contient \emph{toutes} les aides de la Région (sport, culture,
entreprises\dots), pas seulement celles qui concernent la recherche. Le choix fait est de
tout collecter : trier par pertinence est le rôle de l'étape de matching, pas du scraper.

% ===========================================================================
\section{Étape suivante : extraction des champs}
\label{sec:extraction}
% ===========================================================================

Fichier : \code{extraction/field\_extractor.py} (pas encore branché sur \code{main.py}).

\pb Beaucoup de champs importants (date limite, éligibilité, budget, dépenses financées)
sont écrits différemment sur chaque site, souvent seulement dans le PDF. Écrire une
expression régulière par site et par champ ne passe pas à l'échelle.

\sol Utiliser un modèle open source de \emph{question-answering} extractif de type BERT,
qui tourne sur le processeur d'un ordinateur portable. On lui pose une question
(\og Quelle est la date limite de dépôt ? \fg) sur un texte, et il renvoie le passage qui
répond, avec un score de confiance. Les réponses à faible score partent en validation humaine.

Modèles candidats (Hugging Face) :
\begin{itemize}
  \item \code{etalab-ia/camembert-base-squadFR-fquad-piaf} (français) ;
  \item \code{deepset/xlm-roberta-base-squad2} (multilingue, utile pour Horizon en anglais).
\end{itemize}

\impl Déjà en place :
\begin{itemize}
  \item \code{QUESTIONS} : une question par champ à extraire ;
  \item \code{load\_call\_text()} : rassemble le texte de la page et le texte de chaque PDF
        (fichiers \code{.txt}) d'un appel ;
  \item \code{split\_into\_chunks()} : découpe le texte en morceaux d'environ 350 mots, car
        ces modèles lisent au plus 512 \emph{tokens} à la fois ;
  \item \code{extract\_fields()} : renvoie pour l'instant des valeurs vides.
\end{itemize}

Reste à faire : poser chaque question sur chaque morceau, garder la réponse de meilleur
score, et stocker la réponse avec son score. Cela demandera les dépendances
\code{transformers} et \code{torch}.

% ===========================================================================
\section{Ajouter une nouvelle source}
% ===========================================================================

\begin{enumerate}
  \item Créer \code{scrapers/<nom>.py} avec une fonction
        \code{scrape(max\_items=None, with\_pdfs=True)} qui renvoie une liste de lignes créées
        avec \code{new\_record(SOURCE)}.
  \item Toutes les requêtes doivent passer par \code{utils/http\_client.py} (délai et
        \code{robots.txt} gérés automatiquement).
  \item Si la source a une page par appel, réutiliser \code{add\_page\_details()} de
        \code{scrapers/common.py}.
  \item Ajouter une ligne dans le dictionnaire \code{SCRAPERS} de \code{scrapers/\_\_init\_\_.py}.
  \item Ajouter en bas du fichier un bloc \code{if \_\_name\_\_ == "\_\_main\_\_":} pour pouvoir
        tester la source seule avec \code{python -m scrapers.<nom>}.
  \item Lancer \code{python main.py -{}-sources <nom> -{}-max-items 3 -{}-no-pdf} et regarder le
        tableau de remplissage.
\end{enumerate}

Ordre de préférence pour une nouvelle source : API ou open data officiel, puis HTML. Une
API est plus stable, plus rapide et plus respectueuse du site.

% ===========================================================================
\section{Limites connues et corrections}
\label{sec:limites}
% ===========================================================================

\subsection{Corrections apportées le 16/11/2026}

Un premier test de chaque source (\code{-{}-max-items 3 -{}-no-pdf}) a mis en évidence les
problèmes suivants, tous corrigés et décrits dans les sections précédentes.

\begin{center}
\small
\begin{tabularx}{\linewidth}{@{}l X X@{}}
\toprule
\textbf{Source} & \textbf{Problème constaté} & \textbf{Correction} \\
\midrule
Toutes &
Fichiers datés à la minute : deux exécutions dans la même minute écrivaient le même CSV. &
Date à la seconde, commune au CSV et au log. \\
\addlinespace
ANR &
0~\% de dates converties : période coupée sur deux lignes ; \og Fermeture \fg{} absent des
mots-clés ; libellés séparés de leur valeur. &
Lecture du bloc de dates de la carte, puis du calendrier de la page (\code{data-date}) ;
mot-clé ajouté ; libellé recollé à sa valeur. 18 appels sur 24 ont des dates exactes. \\
\addlinespace
ANR &
Textes des appels France 2030 vides (PDF en images, OCR non installé). &
OCR ajouté aux dépendances. \\
\addlinespace
Horizon &
458 topics sur 787 périmés mais marqués ouverts ou à venir. &
Topics dont la dernière date limite est passée ignorés. \\
\addlinespace
Horizon &
Pagination instable : 771 topics différents reçus sur 787. &
Tri par identifiant : 787 sur 787. \\
\addlinespace
Région &
Dates en UTC enregistrées avec un jour de retard ; statut toujours vide ; texte en double ;
pas de PDF. &
Conversion à l'heure de Paris ; statut calculé ; colonnes en double et techniques retirées ;
PDF reliés depuis le second jeu de données. \\
\bottomrule
\end{tabularx}
\end{center}

\subsection{Limites restantes}

\begin{itemize}
  \item \textbf{Bpifrance} : le site renvoie une erreur 403 (protection CloudFront) pour
        toutes les URL, \code{robots.txt} compris. Aucun appel n'est collecté. Ne pas
        contourner la protection ; demander une autorisation ou trouver une source
        alternative.
  \item \textbf{ANR, pré-annonces} : seules des périodes en mois sont publiées ; les champs
        \code{opening\_date} et \code{deadline} restent vides et l'information est dans
        \code{dates\_text}.
  \item \textbf{ANR, sélecteurs CSS} : la lecture des dates repose sur les classes
        \code{date} et \code{task} du site. Si l'ANR refond son site, ces champs se
        videront ; le tableau de remplissage affiché à la fin de chaque exécution permet de
        le voir tout de suite.
  \item \textbf{Région} : pas de type d'appel ; le jeu de données contient toutes les aides
        de la Région, y compris closes (statut \og clos \fg).
  \item \textbf{Horizon} : quelques topics à venir n'ont pas encore de description
        (environ 1~\%).
  \item \textbf{OCR} : il faut installer Tesseract et Poppler sur chaque machine, en plus
        des paquets Python.
  \item \textbf{User-Agent} : l'email de contact de \code{config.py} est encore générique
        et doit être remplacé avant une utilisation régulière.
\end{itemize}

\end{document}
